% Lines 1808-1935 from original attention_transformers_notes_ghost.tex
% Chapter 8: Advanced Topics and Extensions

%==============================================================================
\section{Advanced Topics and Extensions}
%==============================================================================

\subsection{Variants of Attention}

\subsubsection{Local Attention}

For very long sequences, attending to all positions becomes computationally prohibitive. Local attention restricts each position to attend only to a window of positions:

\begin{equation}
\text{LocalAttention}(i) = \text{Attention}(\mathbf{q}_i, \mathbf{K}_{[i-w:i+w]}, \mathbf{V}_{[i-w:i+w]})
\end{equation}
\begin{notationbox}
\textbf{Notation:}
\begin{itemize}
    \item $w$: Window size
    \item $\mathbf{K}_{[i-w:i+w]}$: Keys within the window
    \item Reduces complexity from $O(n^2)$ to $O(n \cdot w)$
\end{itemize}
\end{notationbox}

\subsubsection{Sparse Attention}

Sparse attention patterns reduce computational complexity while maintaining model expressiveness:

\begin{itemize}
    \item \textbf{Strided Attention}: Attend to every $k$-th position
    \item \textbf{Fixed Patterns}: Predefined sparse patterns (e.g., Longformer)
    \item \textbf{Learned Patterns}: Learn which positions to attend to
\end{itemize}

\subsection{Vision Transformer (ViT)}

The Vision Transformer adapts the transformer architecture for image classification:

\begin{figure}[H]
\centering
\begin{tikzpicture}[scale=0.8, every node/.style={scale=0.8}]
    % Original image
    \node[draw, minimum width=3cm, minimum height=3cm, fill=gray!20] (img) at (0,0) {
        \includegraphics[width=2.8cm]{example-image}
    };
    \node[below=0.1cm of img] {Input Image};

    % Patches
    \foreach \i in {0,1,2} {
        \foreach \j in {0,1,2} {
            \pgfmathtruncatemacro\idx{\i*3+\j+1}
            \node[draw, minimum width=0.8cm, minimum height=0.8cm, fill=blue!20]
                (patch\idx) at (5+\j*1,2-\i*1) {\idx};
        }
    }
    \node[below=0.5cm of patch8] {Image Patches};

    % Linear projection
    \node[draw, fill=yellow!20, minimum width=4cm, minimum height=0.8cm]
        (linear) at (9,0) {Linear Projection};

    % Position embeddings
    \node[draw, fill=orange!20, minimum width=4cm, minimum height=0.8cm]
        (pos) at (9,-1.5) {Position Embeddings};

    % Class token
    \node[draw, fill=red!20, minimum width=1cm, minimum height=0.8cm]
        (cls) at (7,-3) {[CLS]};

    % Transformer encoder
    \node[draw, fill=green!20, minimum width=5cm, minimum height=2cm]
        (transformer) at (9,-4.5) {Transformer Encoder};

    % MLP head
    \node[draw, fill=purple!20, minimum width=3cm, minimum height=0.8cm]
        (mlp) at (9,-6.5) {MLP Head};

    % Output
    \node[draw, fill=gray!20, minimum width=2cm, minimum height=0.8cm]
        (output) at (9,-8) {Class};

    % Arrows
    \draw[arrow, thick] (img) -- (patch5);
    \foreach \i in {1,...,9} {
        \draw[arrow] (patch\i) -- (linear);
    }
    \draw[arrow] (linear) -- (pos);
    \draw[arrow] (pos) -- (transformer);
    \draw[arrow] (cls) -- (transformer);
    \draw[arrow] (transformer) -- (mlp);
    \draw[arrow] (mlp) -- (output);

    % Annotations
    \draw[decorate, decoration={brace,amplitude=5pt}]
        (patch1.north west) -- (patch3.north east)
        node[midway,above=5pt] {$P \times P$};
\end{tikzpicture}
\caption{Vision Transformer Architecture}
\label{fig:vit}
\end{figure}

Key differences from NLP transformers:
\begin{enumerate}
    \item \textbf{Input Representation}: Images are divided into fixed-size patches
    \item \textbf{Linear Projection}: Each patch is flattened and linearly projected
    \item \textbf{Classification Token}: A learnable [CLS] token is prepended
    \item \textbf{Position Embeddings}: Learnable 2D position embeddings
    \item \textbf{Output}: Only the [CLS] token representation is used for classification
\end{enumerate}

\subsection{Efficient Transformers}

Recent research has focused on reducing the quadratic complexity of attention:

\begin{table}[H]
\centering
\caption{Efficient Transformer Variants}
\begin{tabular}{l c l}
\toprule
\textbf{Model} & \textbf{Complexity} & \textbf{Method} \\
\midrule
Linformer & $O(n)$ & Low-rank approximation of attention \\
Performer & $O(n)$ & Kernel-based approximation \\
Reformer & $O(n \log n)$ & Locality-sensitive hashing \\
Longformer & $O(n)$ & Sliding window + global attention \\
BigBird & $O(n)$ & Sparse patterns + random attention \\
\bottomrule
\end{tabular}
\end{table}
